\begin{lemma}[No zero divisors in $\mathbb{Q}$]
If $a,b\in\mathbb{Q}$ and $ab=0$, then $a=0$ or $b=0$.
\end{lemma}

\begin{proof}
Suppose $ab=0$.  If $a=0$, there is nothing to prove.  Now assume $a\neq0$.
Then $a^{-1}$ exists in
$\mathbb{Q}$.  Multiplying both sides of $ab=0$ by $a^{-1}$ gives
\[
a^{-1}(ab)=a^{-1}\cdot0.
\]
By associativity, the left side is
\[
(a^{-1}a)b=1b=b.
\]
Also $a^{-1}\cdot0=0$: indeed,
\[
a^{-1}\cdot0=a^{-1}(0+0)=a^{-1}\cdot0+a^{-1}\cdot0,
\]
and adding the additive inverse of $a^{-1}\cdot0$ to both sides gives
$a^{-1}\cdot0=0$.  Hence $b=0$.  Therefore $a=0$ or $b=0$.
\end{proof}